\documentclass[12pt]{article}
\usepackage[a4paper,top=1pt,bottom=2pt,left=1pt,right=1pt,marginparwidth=1pt,headheight=1pt]{geometry}

\input{preamble.tex}

\usepackage{blindtext}
\usepackage{multicol}
\usepackage{color}
\setlength{\columnsep}{0.2cm}
\setlength{\columnseprule}{1pt}
\def\columnseprulecolor{\color{blue}}

\usepackage{xpatch}
\xpatchcmd{\NCC@ignorepar}{%
\abovedisplayskip\abovedisplayshortskip}
{%
\abovedisplayskip\abovedisplayshortskip%
\belowdisplayskip\belowdisplayshortskip}
{}{}

\setlength{\parindent}{0in}
\setlength{\parskip}{0in}

\setlength{\belowdisplayskip}{0pt} \setlength{\belowdisplayshortskip}{0pt}
\setlength{\abovedisplayskip}{0pt} \setlength{\abovedisplayshortskip}{0pt}



\usepackage{soul}
\usepackage[dvipsnames]{xcolor}
\usepackage{listings}
\newcommand{\bulletPoint}[1]{\ul{\textit{\textbf{#1}}}}



% For better text align
\usepackage{ragged2e}
\lstset{
    basicstyle=\ttfamily\tiny,
    breaklines=true,
    breakatwhitespace=false,
    columns=fullflexible,
    keepspaces=true,
    showstringspaces=false
}



\begin{document}
\singlespacing
\begin{multicols*}{4}
\tiny
\raggedright
% Add new cheatsheet content here.

\bulletPoint{Softmax Layer}:
\begin{itemize}
\item In multi-class classification tasks, the output layer is typically a softmax layer, i.e., it employs a softmax activation function.
\item If a layer with a sigmoid activation function is used as the output layer instead, the predictions by the NN may not be easy to interpret, i.e., the outputs do not sum up to 1.
\item Note that an output layer with sigmoid activations can still be used for binary classification.
\item A Softmax layer applies the Softmax function to convert logits into probability values ranging from $[0,1]$.
\item The outputs sum to 1, making it suitable for multi-class classification.
\item The values $z$ inputted to the softmax layer are referred to as logits.
\end{itemize}

\bulletPoint{Activation: Sigmoid}:
\begin{itemize}
\item Sigmoid function $\sigma$: takes a real-valued number and ``squashes'' it into the range between 0 and 1.
\item The output can be interpreted as the firing rate of a biological neuron.
\item Not firing $=0$; Fully firing $=1$.
\item When the neuron's activation are 0 or 1, sigmoid neurons saturate.
\item Gradients at these regions are almost zero, so almost no signal will flow.
\item Sigmoid activations are less common in modern NNs.
\end{itemize}

\bulletPoint{Activation: Tanh}:
\begin{itemize}
\item Tanh function: takes a real-valued number and ``squashes'' it into range between $-1$ and $1$.
\item Like sigmoid, tanh neurons saturate.
\item Unlike sigmoid, the output is zero-centered.
\item It is therefore preferred than sigmoid.
\item Tanh is a scaled sigmoid: $\tanh(x)=2\cdot \sigma(2x)-1$.
\end{itemize}

\bulletPoint{Activation: ReLU}:
\begin{itemize}
\item ReLU (Rectified Linear Unit): takes a real-valued number and thresholds it at zero.
\end{itemize}
\useshortskip
\begin{equation*}
    f(x)=\max(0,x)
\end{equation*}
\begin{itemize}
\item Most modern deep NNs use ReLU activations.
\item ReLU is fast to compute.
\item Compared to sigmoid, tanh, it simply thresholds a matrix at zero.
\item ReLU accelerates the convergence of gradient descent due to its linear, non-saturating form.
\item ReLU prevents the gradient vanishing problem.
\end{itemize}

\bulletPoint{Activation: Leaky ReLU}:
\begin{itemize}
\item The problem of ReLU activations: they can ``die''.
\item ReLU could cause weights to update in a way that the gradients can become zero and the neuron will not activate again on any data.
\item Example: when a large learning rate is used.
\item Leaky ReLU activation function is a variant of ReLU.
\item Instead of the function being 0 when $x<0$, a Leaky ReLU has a small positive slope, e.g., $\alpha=0.01$ or similar.
\item This resolves the dying ReLU problem.
\item Most current works still use ReLU.
\item With a proper setting of the learning rate, the problem of dying ReLU can be avoided.
\end{itemize}

\bulletPoint{Activation: Linear Function}:
\begin{itemize}
\item Linear function means that the output signal is proportional to the input signal to the neuron.
\item If the value of the constant $c$ is 1, it is also called identity activation function.
\item This activation type is used in regression problems.
\item Example: the last layer can have linear activation function in order to output a real number, and not a class membership.
\end{itemize}

\bulletPoint{Activation: Maxout}:
\begin{itemize}
\item Maxout outputs the maximum value among multiple affine functions.
\item For input $\mathbf{x}$, a Maxout unit can be written as:
\end{itemize}
\useshortskip
\begin{equation*}
    f(\mathbf{x})=\max_{i=1,\ldots,k}(\mathbf{w}_i^T\mathbf{x}+b_i)
\end{equation*}
\begin{itemize}
\item It can be viewed as a learnable piecewise linear activation.
\item It does not saturate in the same way as sigmoid/tanh, and it avoids the zero-gradient region of standard ReLU on the negative side.
\item It is more flexible than ReLU/Leaky ReLU because the activation shape is learned from data.
\item Drawback: it introduces more parameters, so computation and memory cost are higher.
\item Typical use case: deep networks that need a more expressive activation, especially when stronger nonlinear modeling is desired.
\item It is also commonly paired with dropout.
\item Example: let $z_1=2x+1$, $z_2=-x+3$. Then
\end{itemize}
\useshortskip
\begin{equation*}
    f(x)=\max(2x+1,-x+3)
\end{equation*}
\begin{itemize}
\item If $x=1$, then $z_1=3$, $z_2=2$, so $f(1)=3$.
\item If $x=-1$, then $z_1=-1$, $z_2=4$, so $f(-1)=4$.
\end{itemize}

\bulletPoint{Activation: ELU}:
\begin{itemize}
\item ELU (Exponential Linear Unit) keeps a non-zero derivative for both positive and negative input, so there is no more dying ReLU problem.
\item Its activation function is
\end{itemize}
\useshortskip
\begin{equation*}
    f(x)=
    \begin{cases}
        x, & x\geq 0\\
        \alpha(e^x-1), & x<0
    \end{cases}
\end{equation*}
\begin{itemize}
\item For $x\ge 0$, ELU behaves like a linear function.
\item For $x<0$, ELU becomes a smooth exponential curve controlled by $\alpha$.
\item Compared with ReLU, ELU allows gradient flow on the negative side.
\end{itemize}

\bulletPoint{Training NNs for Classification}:
\begin{itemize}
\item The network parameters $\theta$ include the weight matrices and bias vectors from all layers.
\item Often, the model parameters $\theta$ are referred to as weights.
\item Training a model means learning a set of parameters $\theta$ that are optimal according to a criterion.
\item This is one of the greatest challenges in ML.
\end{itemize}

\bulletPoint{Data Preprocessing}:
\begin{itemize}
\item Data preprocessing helps convergence during training.
\item Standardization: divide each feature by its standard deviation to obtain standard deviation 1 for each data dimension (feature).
\item Normalization (by mean subtraction / scaling): scale the data within $[0,1]$ or $[-1,1]$.
\item Example: image pixel intensities are divided by 255 to be scaled into the $[0,1]$ range.
\item PyTorch example:
\end{itemize}
\useshortskip
\begin{equation*}
    x_{\text{std}}=\frac{x-\mu}{\sigma}
\end{equation*}
\useshortskip
\begin{equation*}
    x_{\text{norm}}=\frac{x-x_{\min}}{x_{\max}-x_{\min}}
\end{equation*}
\texttt{import torch}\\
\texttt{x = torch.randn([100])}\\
\texttt{x.mean() = 0.02, x.var() = 0.9}\\
\texttt{x\_standardized = (x - x.mean()) / x.std()}\\
\texttt{x\_standardized.mean() = 0.0, x\_standardized.var() = 1.0}\\
\texttt{x\_normalized = (x - x.min()) / (x.max() - x.min())}\\
\texttt{x\_normalized.min() = 0.0, x\_normalized.max() = 1.0}

\bulletPoint{Cross-Entropy Example}:
\texttt{import torch}\\
\texttt{from torch.nn.functional import cross\_entropy}\\
\texttt{y = torch.tensor([0.8, 0.1, 0.06, 0.04])}\\
\texttt{cross\_entropy(y, torch.tensor(0)) = 0.8926}\\
\texttt{cross\_entropy(y, torch.tensor(1)) = 1.5926}\\
\texttt{cross\_entropy(y, torch.tensor(2)) = 1.6326}\\
Predicted probabilities: $P(y=0)=0.8,\; P(y=1)=0.1,\; P(y=2)=0.06,\; P(y=3)=0.04$.\\
Smaller CE loss means the predicted distribution matches the true label better.

\bulletPoint{Gradient Descent Algorithm}:
\begin{itemize}
\item Gradient descent algorithm stops when a local minimum of the loss surface is reached.
\item GD does not guarantee reaching a global minimum.
\item However, empirical evidence suggests that GD works well for NNs.
\item For most tasks, the loss surface $\mathcal{L}(\theta)$ is highly complex and non-convex.
\item Random initialization in NNs results in different initial parameters $\theta_0$ every time the NN is trained.
\item Gradient descent may reach different minima at every run.
\item Therefore, a NN may produce different predicted outputs.
\item Currently, we do not have algorithms that guarantee reaching a global minimum for an arbitrary loss function.
\item Besides the local minima problem, GD can be very slow at plateaus, and it can get stuck at saddle points.
\end{itemize}

\bulletPoint{Mini-Batch Gradient Descent}:
\begin{itemize}
\item Therefore, GD, also known as vanilla GD, is almost always replaced with mini-batch GD.
\item Approach: compute the loss $\mathcal{L}(\theta)$ on a mini-batch of images, update the parameters $\theta$, and repeat until all images are used.
\item At the next epoch, shuffle the training data and repeat the above process.
\item It works because the gradient from a mini-batch is a good approximation of the gradient from the entire training set.
\end{itemize}

\bulletPoint{Adam}:
\begin{itemize}
\item Adam stands for Adaptive Moment Estimation.
\item Adam computes a weighted average of past gradients, i.e., the first moment of the gradient.
\item Adam also computes a weighted average of past squared gradients, i.e., the second moment of the gradient.
\item Other commonly used optimization methods include Adagrad, Adadelta, RMSprop, Nadam, etc.
\item The most commonly used optimizers nowadays are Adam and SGD with momentum.
\end{itemize}

\bulletPoint{Learning Rate}:
\begin{itemize}
\item The gradient tells us the direction in which the loss has the steepest rate of increase.
\item However, it does not tell us how far along the opposite direction we should step.
\end{itemize}

\bulletPoint{Learning Rate Scheduling}:
\begin{itemize}
\item Learning rate scheduling changes the value of the learning rate during training.
\item Annealing means reducing the learning rate over time, also called learning rate decay.
\item Approach 1: reduce the learning rate by some factor every few epochs.
\item Typical values: reduce the learning rate by half every 5 epochs, or divide by 10 every 20 epochs.
\item Approach 2: exponential decay or cosine decay gradually reduce the learning rate over time.
\item Approach 3: reduce the learning rate by a constant, e.g., by half, whenever the validation loss stops improving.
\item This requires splitting data into a training set and a validation set.
\item In TensorFlow: \texttt{tf.keras.callbacks.ReduceLROnPlateau()}.
\item Common settings: monitor validation loss, factor $=0.1$ (divide by 10), patience $=10$, minimum learning rate $=10^{-6}$.
\item Warmup means gradually increasing the learning rate initially, and afterward letting it cool down until the end of training.
\item Barebones PyTorch backpropagation example:
\end{itemize}
\begin{lstlisting}
optimizer = torch.optim.SGD(network.parameters(), lr=0.1)
scheduler = torch.optim.lr_scheduler.ExponentialLR(optimizer, gamma=0.9)

for epoch in range(epochs):
    for x, y in zip(X, Y):
        optimizer.zero_grad()   # Reset gradients
        y_pred = network(x)
        loss = criterion(y_pred, y)
        loss.backward()         # Backpropagate, calculate gradients
        optimizer.step()        # Update weights with LR
    scheduler.step()           # Reduce LR for each epoch
\end{lstlisting}
\begin{itemize}
\item Note that the scheduler is per epoch.
\item If it is called in the inner loop, the learning rate will decay much faster because it will be updated once per mini-batch instead of once per epoch.
\end{itemize}

\bulletPoint{Vanishing / Exploding Gradient Problem}:
\begin{itemize}
\item During training, gradients can become either very small (vanishing gradients) or very large (exploding gradients).
\item They result in very small or very large updates of the parameters.
\item Common solutions: ReLU activations and batch normalization.
\end{itemize}

\bulletPoint{Generalization}:
\begin{itemize}
\item Underfitting: the model is too simple to represent all relevant target characteristics.
\item Example: model with too few parameters.
\item Underfitting produces high error on the training set and high error on the validation set.
\item Overfitting: the model learns patterns specific to the training data, including noise, reducing generalization to unseen data.
\item Example: model with too many parameters.
\item Overfitting produces low error on the training set but high error on the validation set.
\end{itemize}

\bulletPoint{$\ell_2$ Regularization}:
\begin{itemize}
\item Add a regularization term that penalizes large weights to the loss function.
\end{itemize}
\useshortskip
\begin{equation*}
    \mathcal{L}_{reg}(\theta)=\mathcal{L}(\theta)+\lambda\sum_k \theta_k^2
\end{equation*}
\begin{itemize}
\item The regularization coefficient $\lambda$ determines how dominant the regularization is during gradient computation.
\item Large $\lambda$ means a stronger penalty on large weights.
\end{itemize}

\bulletPoint{$\ell_1$ Regularization}:
\begin{itemize}
\item The regularization term is based on the $\ell_1$ norm of the weights.
\end{itemize}
\useshortskip
\begin{equation*}
    \mathcal{L}_{reg}(\theta)=\mathcal{L}(\theta)+\lambda\sum_k |\theta_k|
\end{equation*}
\begin{itemize}
\item $\ell_1$ regularization is less common with NNs.
\item It often performs worse than $\ell_2$ regularization.
\item It is also possible to combine $\ell_1$ and $\ell_2$ regularization, called elastic net regularization.
\end{itemize}
\useshortskip
\begin{equation*}
    \mathcal{L}_{reg}(\theta)=\mathcal{L}(\theta)+\lambda_1\sum_k |\theta_k|+\lambda_2\sum_k \theta_k^2
\end{equation*}

\bulletPoint{Regularization: Dropout}:
\begin{itemize}
\item Randomly drop units, together with their connections, during training.
\item Each unit is dropped with a fixed dropout rate $p$, independent of other units.
\item The hyper-parameter $p$ needs to be chosen or tuned.
\item Often, between 20\% and 50\% of the units are dropped.
\item Dropout is a kind of ensemble learning.
\item One mini-batch trains one network with dropout, and each time this results in a different architecture.
\end{itemize}

\bulletPoint{Regularization: Early Stopping}:
\begin{itemize}
\item During model training, use a validation set.
\item Example: validation/train ratio about 25\% to 75\%.
\item Stop when the validation accuracy or validation loss has not improved after $n$ epochs.
\item The parameter $n$ is called patience.
\end{itemize}

\bulletPoint{Batch Normalization}:
\begin{itemize}
\item Batch normalization layers act similarly to data preprocessing.
\item They calculate the mean $\mu$ and variance $\sigma$ of a batch of input data, and normalize the data $x$ to zero mean and unit variance.
\end{itemize}
\useshortskip
\begin{equation*}
    \hat{x}=\frac{x-\mu}{\sigma}
\end{equation*}
\begin{itemize}
\item BatchNorm alleviates the difficulty of proper initialization of parameters and hyper-parameters.
\item It normalizes activations, making them less dependent on the scale of the parameters during initialization.
\item Therefore, it is less likely to suffer from vanishing or exploding gradients during training.
\item BatchNorm reduces internal covariate shift, i.e., the change in the distribution of activations due to parameter updates during training.
\item It stabilizes training and allows higher learning rates without risking divergence.
\item Consequently, models converge faster and more reliably.
\item BatchNorm layers are very common in convolutional NNs.
\item In PyTorch, Batch Normalization is just another layer:
\end{itemize}
\begin{lstlisting}
network = torch.nn.Sequential(
    torch.nn.Linear(2, 2),
    torch.nn.BatchNorm1d(2),
    torch.nn.Linear(2, 2),
)
\end{lstlisting}

\bulletPoint{Hyper-Parameter Tuning}:
\begin{itemize}
\item Training NNs can involve setting many hyper-parameters:
\item Number of layers and number of neurons per layer.
\item Initial learning rate.
\item Learning rate decay schedule, e.g., decay constant.
\item Optimizer type.
\item Regularization parameters, e.g., $\ell_2$ penalty and dropout rate.
\item Batch size.
\item Activation functions.
\item Loss function.
\item Hyper-parameter tuning can be time-consuming for larger NNs.
\item Grid search: check all values in a range with a step value.
\item Random search: randomly sample values for the parameters; often preferred to grid search.
\item Bayesian hyper-parameter optimization is an active area of research.
\end{itemize}

\bulletPoint{k-Fold Cross-Validation}:
\begin{itemize}
\item Common for hyper-parameter tuning when the training data is small.
\item It gives a better and less noisy estimate of model performance by averaging across folds.
\item Example: 5-fold cross-validation.
\item 1. Split the data into training data and test data.
\item 2. Split the training data into 5 equal folds.
\item 3. Use folds 2--5 for training and fold 1 for validation.
\item 4. Repeat with fold 2 for validation, then fold 3, fold 4, and fold 5.
\item 5. Average the results over the 5 runs for reporting.
\item 6. After choosing the best hyper-parameters, evaluate the model on the test data.
\end{itemize}

\bulletPoint{Deep vs Shallow Networks}:
\begin{itemize}
\item Deeper networks often perform better than shallow networks.
\item However, only up to some limit: after a certain number of layers, the performance of deeper networks plateaus.
\end{itemize}

\bulletPoint{Biological vs Artificial NNs}:
\begin{itemize}
\item Biological neurons respond slowly: about $10^{-3}\,$s, whereas artificial circuits operate around $10^{-9}\,$s.
\item The brain uses massively parallel computation.
\item About $10^{11}$ neurons in the brain.
\item About $10^4$ connections per neuron.
\item About $10^{15}$ total connections in the brain.
\item GPT-3 has about $10^{11}$ parameters.
\item GPT-4 has about $10^{14}$ parameters.
\item A typical biological neuron, e.g., a pyramidal cell, has complex dynamics and computations.
\item The brain also uses extremely complex organization.
\item Current NN models mostly deal only with ``neurons'' and ``local circuits''.
\item The human brain is still superior in some computations: speech recognition under noisy conditions, speech production, and vision in complex situations.
\item Recent autonomous driving systems are still not truly fully autonomous.
\item There is still a long way to go before building a computer that truly mimics the human brain.
\item Very little is known about the human brain.
\item Even simulating a single biological neuron can require a supercomputer.
\item It is unclear which brain processes and structures are functionally essential.
\item We may not need to mimic the human brain to achieve superhuman intelligence. Probably not.
\end{itemize}

\bulletPoint{Historical Sketch}:
\begin{itemize}
\item 1940s: McCulloch and Pitts, Hebb.
\item First mathematical neuron model as all-or-none devices.
\item Hebbian learning rule as a mechanism for learning in biological neurons.
\item 1950s: Rosenblatt, single-layer perceptrons.
\item 1960s: Minsky and Papert published the limitations of single-layer NNs.
\item Their conjecture that multi-layer networks had the same limitations was wrong.
\item This reduced research funding and suspended much NN research.
\item 1980s: resurgence in the field.
\item Hopfield's spin-like neural networks attracted researchers and funding.
\item Rumelhart et al. re-discovered and popularized backpropagation for multi-layer perceptrons.
\item The basic ideas of BP can be traced to Bryson and Ho (1969), \textit{Adaptive Optimal Control}.
\item 2000s: deep learning is old science with better engineering.
\item Faster computers (GPU, etc.), better software (TensorFlow, etc.), and larger datasets enabled progress.
\item Milestones: AlexNet (2012), AlphaGo defeating Lee Sedol (2016), ChatGPT.
\item Deep learning is used across cloud services, medicine, media, security, and autonomous machines.
\end{itemize}

\bulletPoint{Bias as Additional Input}:
\begin{itemize}
\item The bias $b$ of a neuron can be modeled as an additional input.
\item Let $x_0=+1$ and $w_0=b$.
\end{itemize}
\useshortskip
\begin{equation*}
    v=\sum_{j=0}^{m} w_j x_j
\end{equation*}
\begin{itemize}
\item Then the neuron computes an activation from local field $v$.
\end{itemize}

\bulletPoint{Single-Neuron Perceptron}:
\begin{itemize}
\item For two inputs $\mathbf{p}=[p_1,p_2]^T$, classification is based on
\end{itemize}
\useshortskip
\begin{equation*}
    \mathbf{w}^T\mathbf{p}+b>0
\end{equation*}
\begin{itemize}
\item If true, output is class 1; otherwise class 2 (or 0 with step activation).
\item Whether the threshold is $>0$ or $\ge 0$ usually does not matter in practice.
\end{itemize}

\bulletPoint{Decision Boundary and Region}:
\begin{itemize}
\item Decision boundary:
\end{itemize}
\useshortskip
\begin{equation*}
    \mathbf{w}^T\mathbf{p}+b=0
\end{equation*}
\begin{itemize}
\item Equivalently, $\mathbf{w}^T\mathbf{p}=-b$.
\item All points on the boundary have the same inner product with $\mathbf{w}$.
\item Therefore, the boundary is a line or hyperplane orthogonal to $\mathbf{w}$.
\item The weight vector determines the orientation of the boundary.
\item The weight vector and bias together determine its location.
\item Multiplying both $\mathbf{w}$ and $b$ by the same constant does not change the boundary.
\item Decision region:
\end{itemize}
\useshortskip
\begin{equation*}
    \mathbf{w}^T\mathbf{p}+b>0
\end{equation*}
\begin{itemize}
\item The side pointed to by $\mathbf{w}$ has positive activation and output $+1$.
\item The opposite side has negative activation and output 0 for a step function, or approximately 0 for sigmoid.
\end{itemize}

\bulletPoint{Example: OR}:
\begin{itemize}
\item Training pairs: $(0,0)\to 0$, $(0,1)\to 1$, $(1,0)\to 1$, $(1,1)\to 1$.
\item One valid choice is
\end{itemize}
\useshortskip
\begin{equation*}
    \mathbf{w}=\begin{bmatrix}1\\1\end{bmatrix}
\end{equation*}
\begin{itemize}
\item Pick a boundary point $(0,0.5)$:
\end{itemize}
\useshortskip
\begin{equation*}
    \mathbf{w}^T\mathbf{p}+b=1\cdot 0+1\cdot 0.5+b=0
\end{equation*}
\useshortskip
\begin{equation*}
    b=-0.5
\end{equation*}

\bulletPoint{Example: AND}:
\begin{itemize}
\item Training pairs: $(0,0)\to 0$, $(0,1)\to 0$, $(1,0)\to 0$, $(1,1)\to 1$.
\item Use the same direction
\end{itemize}
\useshortskip
\begin{equation*}
    \mathbf{w}=\begin{bmatrix}1\\1\end{bmatrix}
\end{equation*}
\begin{itemize}
\item Choose boundary point $(1,0.5)$:
\end{itemize}
\useshortskip
\begin{equation*}
    \mathbf{w}^T\mathbf{p}+b=1+0.5+b=0
\end{equation*}
\useshortskip
\begin{equation*}
    b=-1.5
\end{equation*}
\begin{itemize}
\item Same boundary orientation as OR, but shifted by a different bias.
\end{itemize}

\bulletPoint{Perceptron Limitation}:
\begin{itemize}
\item A single-layer perceptron gives only a linear decision boundary.
\item Therefore, it cannot handle linearly inseparable problems such as XOR.
\item Multi-layer perceptrons can handle linearly inseparable problems.
\end{itemize}

\bulletPoint{Example: XOR with Hidden Layers}:
\begin{itemize}
\item XOR data: $(0,0)\to 0$, $(0,1)\to 1$, $(1,0)\to 1$, $(1,1)\to 0$.
\item One construction uses hidden neurons to create intermediate regions, then combines them with AND / OR logic.
\item Example boundaries for a compact XOR solution:
\end{itemize}
\useshortskip
\begin{equation*}
    \begin{split}
        \mathbf{w}_1^T\mathbf{p}-0.5&=0,\quad \mathbf{w}_1=\begin{bmatrix}1\\1\end{bmatrix}\\
        \mathbf{w}_2^T\mathbf{p}+1.5&=0,\quad \mathbf{w}_2=\begin{bmatrix}-1\\-1\end{bmatrix}
    \end{split}
\end{equation*}
\begin{itemize}
\item Output neuron can then perform AND on the hidden outputs, with weights $(1,1)$ and bias $-1.5$.
\end{itemize}

\bulletPoint{3-Layer XOR Network}:
\begin{itemize}
\item Let $\sigma$ be the threshold function.
\end{itemize}
\useshortskip
\begin{equation*}
    \mathbf{h}_1=\sigma(\mathbf{W}_1\mathbf{p}+\mathbf{b}_1)\in\mathbb{R}^4
\end{equation*}
\useshortskip
\begin{equation*}
    \mathbf{h}_2=\sigma(\mathbf{W}_2\mathbf{h}_1+\mathbf{b}_2)\in\mathbb{R}^2
\end{equation*}
\useshortskip
\begin{equation*}
    \hat{y}=\sigma(\mathbf{W}_3\mathbf{h}_2+\mathbf{b}_3)\in\mathbb{R}
\end{equation*}
\begin{itemize}
\item One explicit parameterization:
\end{itemize}
\useshortskip
\begin{equation*}
    \mathbf{W}_1=
    \begin{bmatrix}
    -1 & 0\\
    0 & -1\\
    1 & 0\\
    0 & 1
    \end{bmatrix},\quad
    \mathbf{b}_1=
    \begin{bmatrix}
    0.5\\
    0.75\\
    -1.5\\
    -0.25
    \end{bmatrix}
\end{equation*}
\useshortskip
\begin{equation*}
    \mathbf{W}_2=
    \begin{bmatrix}
    1 & 1 & 0 & 0\\
    0 & 0 & 1 & 1
    \end{bmatrix},\quad
    \mathbf{b}_2=
    \begin{bmatrix}
    -1.5\\
    -1.5
    \end{bmatrix}
\end{equation*}
\useshortskip
\begin{equation*}
    \mathbf{W}_3=
    \begin{bmatrix}
    1 & 1
    \end{bmatrix},\quad
    b_3=-0.5
\end{equation*}
\begin{itemize}
\item Example for $\mathbf{p}=[0,1]^T$:
\end{itemize}
\useshortskip
\begin{equation*}
    \mathbf{h}_1=\sigma\!\left(
    \begin{bmatrix}
    0.5\\-0.25\\-1.5\\0.75
    \end{bmatrix}\right)
    =
    \begin{bmatrix}
    1\\0\\0\\1
    \end{bmatrix}
\end{equation*}
\useshortskip
\begin{equation*}
    \mathbf{h}_2=\sigma\!\left(
    \begin{bmatrix}
    -0.5\\-0.5
    \end{bmatrix}\right)
    =
    \begin{bmatrix}
    0\\0
    \end{bmatrix}
\end{equation*}
\useshortskip
\begin{equation*}
    \hat{y}=\sigma(-0.5)=0
\end{equation*}

\bulletPoint{XOR NumPy Implementation}:
\begin{lstlisting}
import numpy as np

w1 = np.array([[-1, 0], [0, -1], [1, 0], [0, 1]])
w2 = np.array([[1, 1, 0, 0], [0, 0, 1, 1]])
w3 = np.array([[1, 1]])

b1 = np.array([0.5, 0.75, -1.5, -0.25])
b2 = np.array([-1.5, -1.5])
b3 = np.array([-0.5])

p = np.array([0, 1]).T

h1 = (w1 @ p + b1) >= 0
h2 = (w2 @ h1 + b2) >= 0
y = (w3 @ h2 + b3) >= 0
\end{lstlisting}
\begin{itemize}
\item In NumPy, `@` is shorthand for matrix multiplication / dot product.
\end{itemize}

\bulletPoint{Universal Approximation}:
\begin{itemize}
\item A 2-layer network can generate only a single convex decision region.
\item A 3-layer MLP can approximate arbitrary decision regions by combining many local regions.
\item For 2 classes, divide one class region into many small hypercubes.
\item First-layer hidden neurons define the hypercube boundaries.
\item Second-layer hidden neurons perform AND operations.
\item Output neuron performs OR across regions.
\item Universal approximation: the accuracy of an MLP can be made arbitrarily high if there are enough hidden neurons.
\item Function approximation theorem: a hidden layer of nonconstant, bounded, continuous neurons plus a linear output neuron can approximate any continuous function to arbitrary accuracy, given enough hidden neurons.
\end{itemize}

\bulletPoint{PyTorch XOR Training Example}:
\texttt{import torch}\\
\texttt{\# XOR dataset}\\
\texttt{X = torch.tensor([[0,0],[0,1],[1,0],[1,1]], dtype=torch.float)}\\
\texttt{Y = torch.tensor([0,1,1,0], dtype=torch.long)}\\
\texttt{network = torch.nn.Linear(2, 2)}\\
\texttt{loss\_fn = torch.nn.CrossEntropyLoss()}\\
\texttt{optim = torch.optim.SGD(network.parameters(), lr=0.1)}\\
\texttt{epochs = 50}\\
\texttt{for epoch in range(epochs):}\\
\hspace*{1em}\texttt{for x, y in zip(X, Y):}\\
\hspace*{2em}\texttt{optim.zero\_grad()}\\
\hspace*{2em}\texttt{y\_pred = network(x)}\\
\hspace*{2em}\texttt{loss = loss\_fn(y\_pred, y)}\\
\hspace*{2em}\texttt{loss.backward()}\\
\hspace*{2em}\texttt{optim.step()}

\bulletPoint{BP vs BP + Momentum Example}:
\begin{lstlisting}
# Sample Python Program for a 1-2-1 neural network, trained with BP vs BP + momentum
# The input data consists of two samples, [0] and [1]
# The corresponding output data is the same as the input data, [0] and [1]
# Prints the loss (mean squared error) for every 1000 epochs
import numpy as np

class NeuralNetwork:
    def __init__(self, input_size, hidden_size, output_size, learning_rate, momentum=None):
        self.input_size = input_size
        self.hidden_size = hidden_size
        self.output_size = output_size
        self.learning_rate = learning_rate
        self.momentum = momentum

        # Initialize weights and biases randomly
        self.weights_input_hidden = np.random.randn(self.input_size, self.hidden_size)
        self.bias_hidden = np.zeros((1, self.hidden_size))
        self.weights_hidden_output = np.random.randn(self.hidden_size, self.output_size)
        self.bias_output = np.zeros((1, self.output_size))

        # Initialize velocities for momentum
        if self.momentum:
            self.velocity_input_hidden = np.zeros((self.input_size, self.hidden_size))
            self.velocity_hidden_output = np.zeros((self.hidden_size, self.output_size))

    def sigmoid(self, x):
        return 1 / (1 + np.exp(-x))

    def sigmoid_derivative(self, x):
        return x * (1 - x)

    def forward_pass(self, X):
        # Forward pass through the network
        self.hidden_output = self.sigmoid(np.dot(X, self.weights_input_hidden) + self.bias_hidden)
        self.output = self.sigmoid(np.dot(self.hidden_output, self.weights_hidden_output) + self.bias_output)
        return self.output

    def backward_pass(self, X, y, output):
        # Backward pass through the network
        error = y - output
        output_delta = error * self.sigmoid_derivative(output)  # *: element-wise multiplication

        # matrix multiplication between output_delta and the transpose of self.weights_hidden_output
        error_hidden = output_delta.dot(self.weights_hidden_output.T)
        hidden_delta = error_hidden * self.sigmoid_derivative(self.hidden_output)

        # Update weights and biases
        if self.momentum:
            self.velocity_input_hidden = self.momentum * self.velocity_input_hidden + X.T.dot(hidden_delta) * self.learning_rate
            self.velocity_hidden_output = self.momentum * self.velocity_hidden_output + self.hidden_output.T.dot(output_delta) * self.learning_rate
            self.weights_input_hidden += self.velocity_input_hidden
            self.weights_hidden_output += self.velocity_hidden_output
        else:
            self.weights_hidden_output += self.hidden_output.T.dot(output_delta) * self.learning_rate
            self.weights_input_hidden += X.T.dot(hidden_delta) * self.learning_rate

        self.bias_output += np.sum(output_delta, axis=0, keepdims=True) * self.learning_rate
        self.bias_hidden += np.sum(hidden_delta, axis=0, keepdims=True) * self.learning_rate

    def train(self, X, y, epochs):
        for epoch in range(epochs):
            # Forward pass
            output = self.forward_pass(X)

            # Backward pass
            self.backward_pass(X, y, output)

            # Print loss for every 1000 epochs
            if epoch % 1000 == 0:
                loss = np.mean(np.square(y - output))
                print(f'Epoch {epoch}, Loss: {loss:.4f}')

# Example usage without momentum
if __name__ == "__main__":
    # Define input data and labels
    X = np.array([[0], [1]])
    y = np.array([[0], [1]])

    # Initialize and train the neural network without momentum
    nn_no_momentum = NeuralNetwork(input_size=1, hidden_size=2, output_size=1, learning_rate=0.1)
    print("Training without momentum:")
    nn_no_momentum.train(X, y, epochs=10000)

    # Initialize and train the neural network with momentum
    nn_with_momentum = NeuralNetwork(input_size=1, hidden_size=2, output_size=1, learning_rate=0.1, momentum=0.9)
    print("\nTraining with momentum:")
    nn_with_momentum.train(X, y, epochs=10000)
\end{lstlisting}


\end{multicols*}
\end{document}
